\chapter{Semisimple Algebras and Induction}

\begin{center}\itshape
The group algebra $\C[G]$ is a product of matrix algebras, one for each irreducible representation. Induction builds representations of a group from those of its subgroups, and Frobenius reciprocity says it is
adjoint to restriction. With these tools, characters prove two theorems about finite groups for which no other proof was known for decades: Burnside's $p^aq^b$ theorem and Frobenius' theorem on permutation groups.
\end{center}

In 1897 William Burnside, professor of mathematics at the Royal Naval College in Greenwich, published \emph{Theory of Groups of Finite Order}, the first book on the subject in English.
In the preface he explained why it said nothing about groups of matrices: ``while, in the present state of our knowledge, many results in the pure theory are arrived at most readily by dealing with properties of
substitution groups, it would be difficult to find a result that could be most directly obtained by the consideration of groups of linear transformations.''

That same year Frobenius was creating character theory (Volume~2, Chapter~11), and Burnside, who read Frobenius' papers carefully, changed his mind. In 1904 he used characters to prove that a group whose order has only
two prime factors is solvable, a theorem for which a proof without characters was found only around 1970. The second edition of his book, in 1911, was half representation theory, and its preface took the old
sentence back: ``the reason given in the original preface for omitting any account of it no longer holds good. In fact it is now more true to say that for further advances in the abstract theory one must look largely to the
representation of a group as a group of linear substitutions.''

This chapter adds the structural side to Volume~2's character theory. Joseph Wedderburn showed in 1908 that algebras like $\C[G]$ are products of matrix algebras; Frobenius found in 1898 how to induce representations from
subgroups; George Mackey, in the 1950s, gave the rules for restricting induced representations. Then we prove Burnside's theorem, and a theorem of Frobenius about permutation groups whose only known proofs still use
characters.

Representations are on finite-dimensional complex vector spaces, as in Volume~2: a representation of a finite group $G$ is the same as a module over the group algebra $\C[G]$, which has basis the elements of $G$ and
the multiplication extending that of $G$. Rings in this chapter need not be commutative; modules are left modules.

\section{Semisimple modules and Wedderburn's theorem}

\begin{definition}\label{def:semisimple}
A module $M$ over a ring $A$ is \emph{simple} if $M\ne0$ and its only submodules are $0$ and $M$, and \emph{semisimple} if it is a direct sum of simple modules. A ring $A$ is \emph{semisimple} if $A$, as a left module over itself,
is semisimple.
\end{definition}

Maschke's theorem (Volume~2, Theorem~11.3) says exactly that every $\C[G]$-module is semisimple; in particular $\C[G]$ is a semisimple ring. Matrix algebras are semisimple: $M_n(D)$ over a division ring $D$ is the direct sum of its
$n$ columns, each isomorphic to the simple module $D^n$. Schur's lemma (Volume~2, Lemma~11.4) holds in any ring: a homomorphism between simple modules is $0$ or an isomorphism, so the endomorphisms of a simple module form a
division ring; and if the module is a finite-dimensional complex vector space, they are the scalars.

\begin{lemma}\label{lem:endA}
For any ring $A$, the map $A^{\mathrm{op}}\to\operatorname{End}_A(A)$, $a\mapsto(x\mapsto xa)$, is an isomorphism of rings, where $A^{\mathrm{op}}$ is $A$ with the product reversed.
\end{lemma}

\begin{proof}
Write $R_a(x)=xa$. Right multiplications commute with left multiplications, so they are module endomorphisms, and $R_b\circ R_a=R_{ab}$: the product $ab$ of $A$ is the product $b\ast a$ of $A^{\mathrm{op}}$, so
$a\mapsto R_a$ is a ring homomorphism from $A^{\mathrm{op}}$, injective as $R_a(1)=a$. An
endomorphism $\varphi$ is determined by $a=\varphi(1)$, since $\varphi(x)=\varphi(x\cdot1)=x\varphi(1)=xa$.
\end{proof}

\begin{theorem}[Wedderburn, 1908]\label{thm:wedderburn}
A semisimple ring $A$ whose module $A$ is a finite direct sum of simple modules is isomorphic to a product $M_{n_1}(D_1)\times\dots\times M_{n_r}(D_r)$ of matrix rings over division rings. If $A$ is a finite-dimensional algebra over $\C$,
all $D_i=\C$.
\end{theorem}

\begin{proof}
Group the simple summands of $A$ by isomorphism type: $A\cong S_1^{n_1}\oplus\dots\oplus S_r^{n_r}$ with $S_i$ pairwise non-isomorphic. By Schur's lemma, there are no non-zero maps $S_i\to S_j$ for $i\ne j$, so an endomorphism of
$A$ is a block-diagonal family of endomorphisms of the $S_i^{n_i}$, and an endomorphism of $S_i^{n_i}$ is an $n_i\times n_i$ matrix with entries in $D_i'=\operatorname{End}_A(S_i)$, a division ring. Hence, with Lemma~\ref{lem:endA},
\[
A^{\mathrm{op}}\cong\operatorname{End}_A(A)\cong\prod_iM_{n_i}(D_i').
\]
Taking opposite rings, and using that the transpose is an isomorphism $M_n(D)^{\mathrm{op}}\cong M_n(D^{\mathrm{op}})$, gives $A\cong\prod M_{n_i}(D_i)$ with $D_i=D_i'^{\mathrm{op}}$. Over $\C$, each $S_i$ is a finite-dimensional
complex vector space, and $D_i=\C$ by Schur's lemma.
\end{proof}

\begin{corollary}\label{cor:wedderburnCG}
Let $\rho_1,\dots,\rho_r$ be the irreducible representations of $G$, of dimensions $d_1,\dots,d_r$. The map
\[
\C[G]\to\prod_{i=1}^r\operatorname{End}(V_i)\cong\prod_{i=1}^rM_{d_i}(\C),\qquad a\mapsto\bigl(\rho_1(a),\dots,\rho_r(a)\bigr),
\]
is an isomorphism of algebras. The centre of $\C[G]$, spanned by the class sums $\sum_{g\in C}g$, corresponds to the scalar matrices, and the element
\[
e_i=\frac{d_i}{|G|}\sum_{g\in G}\chi_i(g^{-1})\,g
\]
is the central idempotent that is the identity on $V_i$ and $0$ on $V_j$ for $j\ne i$.
\end{corollary}

\begin{proof}
The map is an algebra homomorphism. If $\rho_i(a)=0$ for all $i$, then $a$ acts as $0$ on every representation (Maschke), in particular on $\C[G]$, and $a=a\cdot1=0$. So it is injective, and by dimensions it is bijective:
$\sum_id_i^2=|G|$ (Volume~2, Corollary~11.9). The centre of a product of matrix algebras is the product of the scalars, of dimension $r$, which is the number of conjugacy classes (Volume~2, Theorem~11.10); the class sums are central
and linearly independent, so they span the centre. Finally $e_i$ is central (its coefficients form a class function), so by Schur's lemma $\rho_j(e_i)$ is a scalar $\lambda$, and taking traces,
\[
\lambda d_j=\frac{d_i}{|G|}\sum_g\chi_i(g^{-1})\chi_j(g)=d_i\langle\chi_j,\chi_i\rangle,
\]
which is $d_i$ if $j=i$ and $0$ otherwise, by orthogonality (Volume~2, Theorem~11.8).
\end{proof}

\begin{figure}[ht]
\centering
\fig{ch08-wedderburn}
\caption{Wedderburn's theorem for $\mathfrak S_4$: under the isomorphism of Corollary~\ref{cor:wedderburnCG}, an element of $\C[\mathfrak S_4]$ is a block-diagonal matrix with one block for each irreducible character (Volume~2, Example~11.13), and the
$24$ coordinates of $\C[\mathfrak S_4]$ become the $1+1+4+9+9$ free entries of the blocks.}
\label{fig:wedderburn}
\end{figure}

\begin{example}\label{ex:wedderburnsmall}
$\C[\mathfrak S_3]\cong\C\times\C\times M_2(\C)$, matching $6=1+1+4$, and $\C[\mathfrak S_4]$ is drawn in Figure~\ref{fig:wedderburn}. Over $\R$ the theorem gives products of matrix algebras over $\R$, $\C$ and Hamilton's quaternions $\mathbb H$, the only finite-dimensional real
division algebras (Frobenius, 1878): for instance $\R[\Z/3]\cong\R\times\C$, and $\R[Q_8]\cong\R^4\times\mathbb H$ for the quaternion group (Exercise~\ref{ex:rq8}). The quaternions of Volume~2, Chapter~12, are a factor of a group
algebra.
\end{example}

\section{Induced representations}

\begin{definition}\label{def:induced}
Let $H$ be a subgroup of $G$ and $W$ a representation of $H$. The \emph{induced representation} is
\[
\operatorname{Ind}_H^GW=\{f\colon G\to W:\ f(hx)=h\cdot f(x)\text{ for all }h\in H,\ x\in G\},
\]
with $G$ acting by right translation, $(g\cdot f)(x)=f(xg)$. For a representation $V$ of $G$, $\operatorname{Res}_HV$ is $V$ viewed as a representation of $H$.
\end{definition}

A function in $\operatorname{Ind}_H^GW$ is determined by its values on representatives $x_1,\dots,x_m$ of the right cosets $Hx_i$, $m=[G:H]$, and these values are arbitrary, so $\dim\operatorname{Ind}_H^GW=[G:H]\dim W$.
For $W$ trivial, $\operatorname{Ind}_H^G\C$ is the space of functions on $H\backslash G$, the permutation representation on the cosets; for $H=\{1\}$ it is the regular representation.

\begin{theorem}[Frobenius reciprocity, 1898]\label{thm:frobrecip}
For representations $V$ of $G$ and $W$ of $H$, the map $T\mapsto\mathrm{ev}_1\circ T$, where $\mathrm{ev}_1(f)=f(1)$, is an isomorphism
\[
\operatorname{Hom}_G\bigl(V,\operatorname{Ind}_H^GW\bigr)\cong\operatorname{Hom}_H\bigl(\operatorname{Res}_HV,W\bigr).
\]
So induction is right adjoint to restriction (Chapter~1), and $\langle\chi_V,\operatorname{Ind}_H^G\psi\rangle_G=\langle\operatorname{Res}_H\chi_V,\psi\rangle_H$ for the characters.
\end{theorem}

\begin{proof}
If $T$ is $G$-linear, $S(v)=T(v)(1)$ is $H$-linear: $S(hv)=(h\cdot Tv)(1)=(Tv)(h)=h\cdot(Tv)(1)=h\cdot S(v)$. Conversely, given an $H$-linear $S$, put $T(v)(x)=S(xv)$. Then $T(v)\in\operatorname{Ind}_H^GW$, since
$T(v)(hx)=S(hxv)=h\cdot S(xv)$, and $T$ is $G$-linear, since $T(gv)(x)=S(xgv)=T(v)(xg)=(g\cdot T(v))(x)$. The two constructions are inverse: $T(v)(1)=S(v)$, and if $S=\mathrm{ev}_1\circ T$ then
$S(xv)=T(xv)(1)=(x\cdot Tv)(1)=(Tv)(x)$. For characters, the dimension of $\operatorname{Hom}_G(V,V')$ is $\langle\chi_V,\chi_{V'}\rangle$ (by Volume~2, Corollary~11.9, and Schur's lemma, after decomposing both sides), and likewise
for $H$.
\end{proof}

\begin{proposition}[the induced character]\label{prop:indchar}
If $\psi$ is the character of $W$, the character of $\operatorname{Ind}_H^GW$ is
\[
\operatorname{Ind}_H^G\psi(g)=\frac1{|H|}\sum_{\substack{x\in G\\ xgx^{-1}\in H}}\psi(xgx^{-1}).
\]
\end{proposition}

\begin{proof}
Identify $\operatorname{Ind}_H^GW$ with $W^m$ by $f\mapsto(f(x_1),\dots,f(x_m))$. For each $i$, write $x_ig=h_ix_{j}$ with $h_i\in H$ and $j=j(i)$. Then $(g\cdot f)(x_i)=f(x_ig)=h_i\cdot f(x_j)$: the matrix of $g$ is a block matrix whose
$i$-th block row has the single block $\rho_W(h_i)$ in column $j(i)$. Its trace is the sum of $\psi(h_i)$ over the $i$ with $j(i)=i$, that is, with $x_igx_i^{-1}=h_i\in H$. Each such coset $Hx_i$ contributes the same value for all
its $|H|$ elements $x=hx_i$, as $\psi(hx_igx_i^{-1}h^{-1})=\psi(x_igx_i^{-1})$; hence the formula.
\end{proof}

\begin{example}\label{ex:dihedral}
Let $G=D_n$ be the dihedral group of order $2n$, with rotations $r^k$ and reflections, and $H=\langle r\rangle\cong\Z/n$. Induce the character $\psi_j(r^k)=\zeta^{jk}$, $\zeta=e^{2\pi i/n}$. A reflection $s$ is conjugate into $H$ by no
element, so $\operatorname{Ind}\psi_j(s)=0$; and $\operatorname{Ind}\psi_j(r^k)=\zeta^{jk}+\zeta^{-jk}=2\cos(2\pi jk/n)$, since $r^k$ is conjugate in $H$ to itself and by a reflection to $r^{-k}$. For $0<j<n/2$ these are the
$2$-dimensional irreducible representations of $D_n$, realised by rotation by $2\pi j/n$ and a reflection of the plane (Exercise~\ref{ex:dihedralmackey}).
\end{example}

\section{Mackey's formula}

For subgroups $H,K$ of $G$, $G$ is the disjoint union of the \emph{double cosets} $HsK$. For $s\in G$, put $K_s=K\cap s^{-1}Hs$, and let $W^{(s)}$ be the representation of $K_s$ on $W$ given by $k\mapsto\rho_W(sks^{-1})$, which makes
sense because $sks^{-1}\in H$.

\begin{theorem}[Mackey, 1951]\label{thm:mackey}
For a set of representatives $s$ of the double cosets $HsK$,
\[
\operatorname{Res}_K\operatorname{Ind}_H^GW\cong\bigoplus_{s}\operatorname{Ind}_{K_s}^KW^{(s)}.
\]
\end{theorem}

\begin{proof}
A function $f\in\operatorname{Ind}_H^GW$ is the sum of its restrictions to the double cosets $D=HsK$, each extended by $0$, and each such restriction is again in $\operatorname{Ind}_H^GW$ (as $HD=D$). Right translation by $K$
preserves each $D$ (as $DK=D$), so $\operatorname{Res}_K\operatorname{Ind}_H^GW=\bigoplus_DF_D$, with $F_D$ the functions supported on $D$. For $f\in F_{HsK}$ put $\varphi(k)=f(sk)$. For $k_0\in K_s$,
\[
\varphi(k_0k)=f(sk_0k)=f\bigl((sk_0s^{-1})sk\bigr)=sk_0s^{-1}\cdot f(sk)=k_0\cdot_{(s)}\varphi(k),
\]
so $\varphi\in\operatorname{Ind}_{K_s}^KW^{(s)}$; and $f\mapsto\varphi$ commutes with right translation by $K$. It is injective, since $f(hsk)=h\cdot f(sk)$. It is onto: given $\varphi$, set $f(hsk)=h\cdot\varphi(k)$. This is well defined, because
$hsk=h'sk'$ gives $k'k^{-1}=s^{-1}h'^{-1}hs\in K_s$, and then $h'\cdot\varphi(k')=h'\cdot\bigl(s(k'k^{-1})s^{-1}\bigr)\cdot\varphi(k)=h'h'^{-1}h\cdot\varphi(k)=h\cdot\varphi(k)$.
\end{proof}

\begin{corollary}[Mackey's irreducibility criterion]\label{cor:mackeyirr}
$\operatorname{Ind}_H^GW$ is irreducible if and only if $W$ is irreducible and, for every $s\in G\setminus H$, the representations $W^{(s)}$ and $\operatorname{Res}_{H_s}W$ of $H_s=H\cap s^{-1}Hs$ have no irreducible constituent in common.
\end{corollary}

\begin{proof}
By Frobenius reciprocity, then Mackey's formula with $K=H$, then Frobenius reciprocity again,
\[
\langle\operatorname{Ind}W,\operatorname{Ind}W\rangle_G=\langle\operatorname{Res}_H\operatorname{Ind}W,W\rangle_H=\sum_s\langle\operatorname{Ind}_{H_s}^HW^{(s)},W\rangle_H=\sum_s\langle W^{(s)},\operatorname{Res}_{H_s}W\rangle_{H_s}.
\]
The term $s=1$ (the double coset $H$) is $\langle W,W\rangle_H\ge1$, with equality if and only if $W$ is irreducible; the other terms are $\ge0$. The total is $1$, that is, $\operatorname{Ind}W$ is irreducible (Volume~2, Corollary~11.9(b)), if and only
if the first term is $1$ and all others vanish.
\end{proof}

\begin{think}
Take $W$ trivial. Then $\langle\operatorname{Ind}_H^G1,\operatorname{Ind}_K^G1\rangle$ counts double cosets $H\backslash G/K$ (Exercise~\ref{ex:doublecosets}). What does this say about the number of orbits of $G$ on pairs of points of $G/H$, and when is
the permutation representation on $G/H$ the trivial representation plus one irreducible?
\end{think}

\section{Burnside's \texorpdfstring{$p^aq^b$}{paqb} theorem}

Characters take values that are algebraic integers, roots of monic polynomials with integer coefficients. By Corollary~\ref{cor:integralclosure}, applied to $\Z\subseteq\C$, the algebraic integers form a ring, and
a rational algebraic integer is an integer ($\Z$ is a principal ideal domain, hence integrally closed, as in the proof of Theorem~\ref{thm:dvr}, (a)$\Rightarrow$(b)).

\begin{lemma}\label{lem:omega}
Let $\chi$ be an irreducible character of degree $d$, and $C$ the conjugacy class of $g$. Then $\chi(g)$ and $\omega=|C|\chi(g)/d$ are algebraic integers.
\end{lemma}

\begin{proof}
$\chi(g)$ is a sum of eigenvalues of $\rho(g)$, which are roots of unity, hence roots of $x^m-1$. For $\omega$: the class sum $c_C=\sum_{x\in C}x$ is central in $\C[G]$, so $\rho(c_C)$ is a scalar (Schur), and taking traces it is
$|C|\chi(g)/d=\omega_C$. The product $c_Cc_D$ is a combination $\sum_Ea_{CDE}c_E$ with non-negative integers $a_{CDE}$ (it is central with integer coefficients, constant on classes). Applying $\rho$, $\omega_C\omega_D=\sum_Ea_{CDE}\omega_E$. So $M=\sum_E\Z\omega_E$ is a subring of $\C$ (it contains $\omega_{\{1\}}=1$ and is closed under products), finitely generated as a $\Z$-module, and it contains $\omega_C$.
By Proposition~\ref{prop:integralcrit}(c), $\omega_C$ is integral over $\Z$.
\end{proof}

\begin{lemma}[Burnside]\label{lem:burnsidevanish}
If, moreover, $|C|$ is prime to $d$, then either $\chi(g)=0$ or $\rho(g)$ is a scalar matrix.
\end{lemma}

\begin{proof}
Write $u|C|+vd=1$ with $u,v\in\Z$. Then $\alpha=\chi(g)/d=u\omega+v\chi(g)$ is an algebraic integer, and $|\alpha|\le1$, since $\chi(g)$ is a sum of $d$ roots of unity $\lambda_1,\dots,\lambda_d$. If $|\alpha|=1$, the triangle inequality
is an equality, so all $\lambda_k$ are equal and $\rho(g)$ (diagonalisable, of finite order) is the scalar $\lambda_1$. Suppose $|\alpha|<1$. Let $m$ be the order of $g$; then $\alpha\in L=\Q(\zeta_m)$, a Galois extension of $\Q$
(Volume~2, Proposition~10.9). For $\sigma\in\operatorname{Gal}(L/\Q)$, $\sigma(\alpha)=\frac1d\sum\sigma(\lambda_k)$ is again an average of roots of unity, so $|\sigma(\alpha)|\le1$, and it is an algebraic integer (it satisfies the same monic
equation as $\alpha$). The product $\beta=\prod_\sigma\sigma(\alpha)$ is an algebraic integer fixed by the Galois group, hence rational (Volume~2, Theorem~10.6), hence an integer; and $|\beta|\le|\alpha|<1$. So $\beta=0$, some
$\sigma(\alpha)=0$, and $\alpha=0$.
\end{proof}

\begin{proposition}\label{prop:primepowerclass}
If a finite group $G$ has a conjugacy class of size $p^k$ with $p$ prime and $k\ge1$, then $G$ is not simple.
\end{proposition}

\begin{proof}
Let $g\in C$, $|C|=p^k$, so $g\ne1$. Suppose $G$ simple; it is non-abelian since $|C|>1$. Column orthogonality (Volume~2, Corollary~11.11) for $g$ and $1$ gives $\sum_\chi\chi(1)\chi(g)=0$, the sum over the irreducible
characters; the trivial character contributes $1$. If every non-trivial $\chi$ with $\chi(g)\ne0$ had $p\mid\chi(1)$, we would get $1+p\gamma=0$ with $\gamma=\sum\frac{\chi(1)}p\chi(g)$ an algebraic integer, and $-\frac1p=\gamma$
would be a rational algebraic integer, which is absurd. So some non-trivial $\chi$ has $p\nmid\chi(1)$ and $\chi(g)\ne0$. By Lemma~\ref{lem:burnsidevanish}, $\rho_\chi(g)$ is a scalar, so it commutes with $\rho_\chi(G)$. The kernel of $\rho_\chi$ is
a normal subgroup, not $G$ as $\chi$ is non-trivial, hence trivial; so $\rho_\chi$ is injective, $g$ commutes with every element of $G$, and $|C|=1$, a contradiction.
\end{proof}

\begin{theorem}[Burnside, 1904]\label{thm:burnsidepq}
Every group of order $p^aq^b$, with $p,q$ primes, is solvable.
\end{theorem}

\begin{proof}
By induction on $|G|$. By Volume~2, Proposition~5.10(b), it suffices to show that $G$ is abelian or has a normal subgroup $N\ne1,G$: then $N$ and $G/N$, of smaller orders of the same form, are solvable. If $b=0$, $G$ is a
$p$-group, solvable by Volume~2, Proposition~5.10(c). Otherwise let $Q$ be a Sylow $q$-subgroup ($b\ge1$) and $1\ne g\in Z(Q)$ (Volume~2, Corollary~3.10). The centraliser of $g$ contains $Q$, so the class of $g$ has size
$[G:C_G(g)]$ dividing $[G:Q]=p^a$. If this size is $1$, $g$ is central and $\langle g\rangle$ is a normal subgroup, which is $\ne G$ unless $G$ is cyclic. Otherwise Proposition~\ref{prop:primepowerclass} gives a normal subgroup
$N\ne1,G$.
\end{proof}

The theorem is sharp in a famous way: the smallest non-abelian simple group, $A_5$, has order $60=2^2\cdot3\cdot5$, with three primes. A proof avoiding characters was found by David Goldschmidt (for $p,q$ odd, 1970), Helmut Bender (1972) and
Hiroshi Matsuyama (1973); it is several times longer.

\section{Frobenius groups}

Let $G$ act transitively on a finite set $X$ with $|X|>1$, so that no element other than $1$ fixes two points. Let $H$ be the stabiliser of a point; the hypothesis says $H\cap xHx^{-1}=\{1\}$ for $x\notin H$. Such a $G$ is a
\emph{Frobenius group}, with \emph{complement} $H$. Example: the affine maps $x\mapsto ax+b$ of a finite field $\mathbb F_q$, with $a\ne0$, acting on $\mathbb F_q$; $H$ is the group of maps $x\mapsto ax$, and the fixed-point-free elements are the
translations $x\mapsto x+b$, $b\ne0$. Here the translations and the identity form a normal subgroup. Frobenius proved in 1901 that this always happens; more than a century later, no proof without characters is known.

\begin{theorem}[Frobenius, 1901]\label{thm:frobeniuskernel}
In a Frobenius group with complement $H$, the set $N$ of elements fixing no point, together with $1$, is a normal subgroup, with $G=NH$ and $N\cap H=\{1\}$.
\end{theorem}

\begin{proof}
\emph{Counting.} The conjugates $xHx^{-1}$ are the stabilisers of the $m=[G:H]=|X|$ points, and two of them meet only in $1$. So their union has $1+m(|H|-1)$ elements, and $N=G\setminus\bigcup_x(xHx^{-1}\setminus\{1\})$ has
$|G|-m(|H|-1)=m$ elements. $N$ is a union of conjugacy classes. If $N$ is a subgroup, it is normal, $N\cap H=\{1\}$, and $|NH|=|N||H|=|G|$.

\emph{Values of induced characters.} Let $\psi$ be a class function on $H$. By Proposition~\ref{prop:indchar}, $\operatorname{Ind}_H^G\psi$ vanishes on $N\setminus\{1\}$ (no conjugate of such an element lies in $H$), equals $m\psi(1)$ at $1$,
and for $1\ne h\in H$ equals $\psi(h)$: if $xhx^{-1}\in H$, then $xhx^{-1}\in H\cap xHx^{-1}$ is not $1$, so $x\in H$, and the $|H|$ such $x$ give the value $\psi(h)$ each.

\emph{An extension of each character.} For an irreducible character $\psi\ne1$ of $H$, let
\[
\theta=\operatorname{Ind}_H^G\psi-\psi(1)\operatorname{Ind}_H^G1_H+\psi(1)1_G .
\]
It is a $\Z$-linear combination of characters. By the previous step, $\theta(1)=m\psi(1)-\psi(1)m+\psi(1)=\psi(1)$, $\theta(h)=\psi(h)-\psi(1)+\psi(1)=\psi(h)$ for $1\ne h\in H$, and $\theta(n)=\psi(1)$ for $n\in N$. As $\theta$ is a class function, $\theta(xhx^{-1})=\psi(h)$ for $1\ne h\in H$, and the $m$ sets
$xHx^{-1}\setminus\{1\}$ are disjoint and cover $G\setminus N$. So
\[
|G|\langle\theta,\theta\rangle=|N|\psi(1)^2+m\sum_{1\ne h\in H}|\psi(h)|^2=m\psi(1)^2+m\bigl(|H|-\psi(1)^2\bigr)=m|H|=|G|,
\]
using $\sum_{h\in H}|\psi(h)|^2=|H|$. Writing $\theta=\sum a_i\chi_i$ with $a_i\in\Z$, $\sum a_i^2=1$, so $\theta=\pm\chi_i$ for one irreducible character, and $\theta(1)>0$ makes it $+\chi_i$: $\theta$ is an irreducible character of $G$, which
restricts to $\psi$ on $H$ and is constant equal to $\theta(1)$ on $N$.

\emph{$N$ is a subgroup.} Let $K$ be the intersection of the kernels of the representations with characters $\theta_\psi$, for all irreducible $\psi\ne1$ of $H$; it is a normal subgroup, since $g\in\ker\rho$ if and only if
$\chi_\rho(g)=\chi_\rho(1)$ (equality in $|\chi(g)|\le\chi(1)$ forces $\rho(g)=1$, as in Lemma~\ref{lem:burnsidevanish}). Every $n\in N$ is in $K$. Conversely, if $g\in K\setminus N$, then $g$ is conjugate to some $1\ne h\in H$, also in $K$, and
$\psi(h)=\theta_\psi(h)=\theta_\psi(1)=\psi(1)$ for every $\psi\ne1$, and also for $\psi=1$. So $h$ acts trivially in every irreducible representation of $H$, hence in the regular one, and $h=1$, a contradiction. So $N=K$.
\end{proof}

\begin{beauty}
The proof never looks inside $N$. It builds, for each irreducible representation of the small group $H$, an irreducible representation of $G$ extending it, and then reads off $N$ as the set of elements that all these
representations ignore. Character-free proofs of Frobenius' theorem are known only in special cases, and the general one remains a well-known open problem. Frobenius also conjectured that $N$ is \emph{nilpotent}; John Thompson
proved it in 1959, in his doctoral thesis, the start of a career that led to the Feit--Thompson theorem of Chapter~8.
\end{beauty}

\section{Exercises}

\begin{exercise}[$\star$]\label{ex:doublecosets}
Show that $\langle\operatorname{Ind}_H^G1,\operatorname{Ind}_K^G1\rangle_G$ is the number of double cosets $HsK$. Deduce that $G$ acts doubly transitively on $X=G/H$ if and only if $\C^X$ is the sum of the trivial representation and one
irreducible representation.
\end{exercise}

\begin{exercise}[$\star$]\label{ex:dihedralmackey}
In Example~\ref{ex:dihedral}, use Mackey's criterion to show that $\operatorname{Ind}\psi_j$ is irreducible if and only if $\psi_j\ne\psi_{-j}$, and find all irreducible representations of $D_n$.
\end{exercise}

\begin{exercise}[$\star\star$]\label{ex:rq8}
Show that $\R[Q_8]\cong\R^4\times\mathbb H$ and $\C[Q_8]\cong\C^4\times M_2(\C)$. (Use the four one-dimensional characters, which factor through $Q_8/\{\pm1\}$, and the action of $Q_8$ on $\mathbb H$ by left multiplication.)
\end{exercise}

\begin{exercise}[$\star\star$]\label{ex:heisenberg}
Let $p$ be an odd prime and $G$ the Heisenberg group of upper unitriangular $3\times3$ matrices over $\mathbb F_p$. Find its $p^2+p-1$ irreducible representations: $p^2$ of dimension $1$, and $p-1$ of dimension $p$, induced from the
abelian subgroup of matrices with zero $(1,2)$ entry.
\end{exercise}

\begin{exercise}[$\star\star$]\label{ex:affinefrob}
For the affine group $G$ of $\mathbb F_q$, show that the characters $\theta_\psi$ of the proof of Theorem~\ref{thm:frobeniuskernel} are the non-trivial characters of $G/N\cong\mathbb F_q^\times$, and that $G$ has exactly one more irreducible
character, of degree $q-1$, induced from a non-trivial character of $N$.
\end{exercise}

\begin{exercise}[$\star\star\star$]\label{ex:centrefaithful}
Show that if $G$ has a faithful irreducible representation, its centre is cyclic, and that the degree of an irreducible representation divides $[G:Z(G)]$. (For the second, apply Lemma~\ref{lem:omega} to the representation of
$G^n$ on $V^{\otimes n}$, as Tate did.)
\end{exercise}

\section{Hints}

\begin{hint}{ex:doublecosets}
Mackey with trivial $W$, or count fixed points: $\operatorname{Ind}1$ is the permutation character.
\end{hint}

\begin{hint}{ex:dihedralmackey}
$H$ is normal, $H_s=H$, and $\psi^{(s)}=\psi_{-j}$ for a reflection $s$.
\end{hint}

\begin{hint}{ex:rq8}
Count dimensions: $8=4\cdot1+4$. The left action of $Q_8$ on $\mathbb H$ commutes with right multiplications by $\mathbb H$.
\end{hint}

\begin{hint}{ex:heisenberg}
The centre $Z$ has order $p$, and $G/Z\cong\mathbb F_p^2$. For a non-trivial character of $Z$, extend it to the abelian subgroup and induce.
\end{hint}

\begin{hint}{ex:affinefrob}
$H=\mathbb F_q^\times$ is abelian; its non-trivial characters have degree $1$.
\end{hint}

\begin{hint}{ex:centrefaithful}
On a faithful irreducible $V$, central elements act by scalars, so $Z(G)$ embeds in $\C^\times$. For the second part, the subgroup $\{(z_1,\dots,z_n):\prod z_i=1\}$ of $Z^n$ acts trivially on $V^{\otimes n}$.
\end{hint}

\section{Solutions}

\begin{solution}{ex:doublecosets}
\idea Frobenius reciprocity and Mackey reduce the inner product to counting.
\writeup By Frobenius reciprocity and Mackey's formula with $W=1$,
\[
\langle\operatorname{Ind}_H1,\operatorname{Ind}_K1\rangle_G=\langle\operatorname{Res}_K\operatorname{Ind}_H1,1\rangle_K=\sum_s\langle\operatorname{Ind}_{K_s}^K1,1\rangle_K=\sum_s\langle1,1\rangle_{K_s},
\]
one term equal to $1$ for each double coset. With $K=H$, $\langle\chi,\chi\rangle$ for the permutation character $\chi$ of $X=G/H$ is the number of double cosets $HsH$, which is the number of orbits of $G$ on $X\times X$ (the orbit of $(H,Hs)$ under
right multiplication determines $HsH$). $G$ is doubly transitive if and only if there are two orbits on $X\times X$ (the diagonal and the rest), that is, $\langle\chi,\chi\rangle=2$. Since $\chi$ contains the trivial character once
($G$ is transitive), $\langle\chi,\chi\rangle=2$ means $\chi=1+\chi'$ with $\chi'$ irreducible.
\end{solution}

\begin{solution}{ex:dihedralmackey}
\idea A normal subgroup of index $2$: one non-trivial double coset, and the criterion compares $\psi_j$ with its conjugate.
\writeup $H=\langle r\rangle$ has index $2$, so the double cosets are $H$ and $sH$ for a reflection $s$, and $H_s=H$. Since $srs^{-1}=r^{-1}$, $\psi_j^{(s)}(r^k)=\psi_j(r^{-k})=\psi_{-j}(r^k)$. By Corollary~\ref{cor:mackeyirr}, $\operatorname{Ind}\psi_j$ is
irreducible iff $\psi_{-j}\ne\psi_j$, i.e. $2j\not\equiv0\pmod n$. Also $\operatorname{Ind}\psi_j\cong\operatorname{Ind}\psi_{-j}$ (same character). This gives $\lfloor(n-1)/2\rfloor$ irreducibles of dimension $2$. The one-dimensional
ones: $D_n/\langle r^2\rangle$ has order $2$ if $n$ is odd and $4$ if $n$ is even, giving $2$ or $4$ characters. Check: for $n$ odd, $2+4\cdot\frac{n-1}2=2n$; for $n$ even, $4+4\cdot\frac{n-2}2=2n$. So these are all.
\end{solution}

\begin{solution}{ex:rq8}
\idea Split $\R[Q_8]$ with the central idempotents $\frac{1\pm z}2$, $z=-1\in Q_8$.
\writeup Let $z$ be the central element $-1$ of $Q_8$ (an element of the group, not the scalar), and $e=\frac{1+z}2$, $e'=\frac{1-z}2$ in $\R[Q_8]$: central idempotents with $e+e'=1$, $ee'=0$. So $\R[Q_8]=e\R[Q_8]\times e'\R[Q_8]$ as
algebras, of dimensions $4$ and $4$. First, $e\R[Q_8]\cong\R[Q_8/\langle z\rangle]=\R[(\Z/2)^2]\cong\R^4$, by the four real characters of $(\Z/2)^2$ (Volume~2, Theorem~11.14). Second, the inclusion $Q_8\subset\mathbb H^\times$ gives an algebra map
$\R[Q_8]\to\mathbb H$ sending $z$ to $-1$, hence $e$ to $0$ and $e'$ to $1$; it is onto, since $1,i,j,k$ are in the image. Restricted to $e'\R[Q_8]$ it is onto $\mathbb H$, and by dimensions an isomorphism. Hence $\R[Q_8]\cong\R^4\times\mathbb H$.
Over $\C$, $\mathbb H\otimes_\R\C\cong M_2(\C)$: the matrices $I=\begin{psmallmatrix}i&0\\0&-i\end{psmallmatrix}$ and $J=\begin{psmallmatrix}0&1\\-1&0\end{psmallmatrix}$ satisfy $I^2=J^2=-1$ and $IJ=-JI$, and $1,I,J,IJ$ span $M_2(\C)$. This matches
$8=4\cdot1^2+2^2$.
\end{solution}

\begin{solution}{ex:heisenberg}
\idea Characters that are trivial on the centre come from $G/Z\cong\mathbb F_p^2$; the others are induced from a maximal abelian subgroup.
\writeup Write $(a,b,c)$ for the matrix whose entries in positions $(1,2)$, $(2,3)$ and $(1,3)$ are $a$, $b$ and $c$; then
\[
(a,b,c)(a',b',c')=(a+a',\,b+b',\,c+c'+ab').
\] The centre is $Z=\{(0,0,c)\}$ and $G/Z\cong\mathbb F_p^2$, giving $p^2$ characters of degree $1$.
The subgroup $A=\{(0,b,c)\}\cong\mathbb F_p^2$ is abelian and normal of index $p$. For $\lambda\ne0$ in $\mathbb F_p$, let $\psi_\lambda(0,b,c)=\zeta^{\lambda c}$, $\zeta=e^{2\pi i/p}$. A representative $s=(a,0,0)$ of $G/A$ conjugates $(0,b,c)$ to $(0,b,c+ab)$, so
$\psi_\lambda^{(s)}(0,b,c)=\zeta^{\lambda(c+ab)}\ne\psi_\lambda$ for $a\ne0$: by Mackey ($A$ normal, $A_s=A$), $\operatorname{Ind}_A^G\psi_\lambda$ is irreducible of dimension $p$, with central character $\zeta^{\lambda c}$; different $\lambda$ give
different central characters, hence non-isomorphic representations. Count: $p^2\cdot1+(p-1)p^2=p^3$, so the list is complete, and there are $p^2+p-1$ irreducibles, the number of conjugacy classes.
\end{solution}

\begin{solution}{ex:affinefrob}
\idea $\theta_\psi$ is $\psi$ on $H$ and $\psi(1)=1$ on $N$, so it factors through $G/N$.
\writeup Here $|N|=q$, $|H|=q-1$ and $m=q$. For a non-trivial character $\psi$ of $H=\mathbb F_q^\times$, of degree $1$, the proof gives $\theta_\psi=1$ on $N$ and $\theta_\psi(x\mapsto ax)=\psi(a)$; as every map $x\mapsto ax+b$ with $a\ne1$ is
conjugate to $x\mapsto ax$ (by a translation), $\theta_\psi(x\mapsto ax+b)=\psi(a)$ for all $a,b$. So $\theta_\psi$ is the character $\psi$ of $G/N\cong\mathbb F_q^\times$. With the trivial character, these are the $q-1$ characters of $G/N$. The
classes of $G$ are the identity, the non-trivial translations (one class: $x\mapsto ax$ conjugates $x+b$ to $x+ab$), and one class for each $a\ne1$: $q$ classes. So one irreducible character is missing, of degree $d$ with
$(q-1)+d^2=q(q-1)$, so $d=q-1$. For a non-trivial character $\lambda$ of $N\cong\mathbb F_q$, the conjugates $b\mapsto\lambda(ab)$, $a\ne1$, differ from $\lambda$, so $\operatorname{Ind}_N^G\lambda$, of degree $q-1$, is irreducible by
Corollary~\ref{cor:mackeyirr}: it is the missing one.
\end{solution}

\begin{solution}{ex:centrefaithful}
\idea Central elements act by scalars; Tate's trick raises the integrality of $\omega$ to a power.
\writeup If $\rho$ is faithful and irreducible, each $z\in Z(G)$ acts by a scalar $\lambda(z)$ (Schur), and $z\mapsto\lambda(z)$ is an injective homomorphism $Z(G)\to\C^\times$; a finite subgroup of $\C^\times$ consists of roots of unity of
some order and is cyclic. For the second part, let $V$ be irreducible of degree $d$, $Z=Z(G)$, and $n\ge1$. $G^n$ acts irreducibly on $V^{\otimes n}$ (its character is $\chi(g_1)\cdots\chi(g_n)$, of norm $1$), with degree $d^n$.
Elements of $Z$ act by scalars $\lambda(z)$, so $(z_1,\dots,z_n)$ acts by $\prod\lambda(z_i)=\lambda(z_1\cdots z_n)$, and the subgroup $Z_0=\{(z_i):\prod z_i=1\}$, of order $|Z|^{n-1}$, acts trivially. So $V^{\otimes n}$ is an irreducible
representation of $G^n/Z_0$, of order $|G|^n/|Z|^{n-1}$, and since degrees divide the group order (Volume~2, Exercise~11.8), $d^n$ divides $|G|^n/|Z|^{n-1}$. With $r=[G:Z]/d\in\Q$, this says that $|Z|\,r^n=|G|^n/(|Z|^{n-1}d^n)$ is an integer for every $n$. If a prime $\ell$ divided the
denominator of $r$ (in lowest terms), $\ell^n$ would divide the denominator of $r^n$, hence $|Z|$, for every $n$: impossible. So $r$ is an integer, that is, $d$ divides $[G:Z]$.
\end{solution}

\section*{Sources and further reading}
\addcontentsline{toc}{section}{Sources and further reading}

W.~Burnside, \emph{Theory of Groups of Finite Order} (Cambridge University Press, 1897; 2nd edition 1911), the prefaces quoted above; W.~Burnside, ``On groups of order $p^\alpha q^\beta$'', \emph{Proc.\ London Math.\ Soc.} (2) 1
(1904), 388--392; G.~Frobenius, ``Über Relationen zwischen den Charakteren einer Gruppe und denen ihrer Untergruppen'', \emph{Sitzungsber.\ Preuss.\ Akad.\ Wiss.} (1898), 501--515, for reciprocity, and ``Über auflösbare Gruppen
IV'', ibid.\ (1901), 1216--1230, for Frobenius groups; J.~H.~M.~Wedderburn, ``On hypercomplex numbers'', \emph{Proc.\ London Math.\ Soc.} (2) 6 (1908), 77--118; G.~W.~Mackey, ``On induced representations of groups'', \emph{Amer.\ J.\ Math.} 73
(1951), 576--592; J.~G.~Thompson, ``Finite groups with fixed-point-free automorphisms of prime order'', \emph{Proc.\ Nat.\ Acad.\ Sci.\ USA} 45 (1959), 578--581.

For more: J.-P.~Serre, \emph{Linear Representations of Finite Groups} (Springer, 1977), Chapters~6 and~7; I.~M.~Isaacs, \emph{Character Theory of Finite Groups} (Academic Press, 1976), Chapters~3, 5 and~7; C.~W.~Curtis,
\emph{Pioneers of Representation Theory: Frobenius, Burnside, Schur, and Brauer} (American Mathematical Society, 1999), for the history.
